% Informe académico - TP Grupal LMC 2C 2026
% Copiar este archivo al proyecto de Overleaf (reemplaza el main.tex actual).
% Los bloques marcados con TODO se completan cuando el notebook esté ejecutado.
% Extensión objetivo: 3 a 5 páginas sin contar portada, índice ni bibliografía.

\documentclass[11pt,a4paper]{article}
\usepackage[utf8]{inputenc}
\usepackage[T1]{fontenc}
\usepackage[spanish]{babel}
\usepackage[margin=2.5cm]{geometry}
\usepackage{graphicx}
\usepackage{booktabs}
\usepackage{float}
\usepackage{amsmath}
\usepackage{hyperref}

\graphicspath{{figuras/}{../graficos/}}

\begin{document}

% ---------------------------------------------------------------------------
% Portada
% ---------------------------------------------------------------------------
\begin{titlepage}
    \centering
    \vspace*{3cm}
    {\LARGE Laboratorio de Métodos Cuantitativos\\ Aplicados a la Gestión\par}
    \vspace{2cm}
    {\Huge\bfseries Análisis F1\par}
    \vspace{0.5cm}
    {\Large Rendimiento por auto o talento\par}
    \vfill
    {\large Alumnos: Agustín Mier Torre y Joaquín Lobato Adinolfi\par}
    \vspace{0.5cm}
    {\large Facultad de Ciencias Económicas, Universidad de Buenos Aires\par}
    \vspace{0.3cm}
    {\large 2do Cuatrimestre 2026\par}
\end{titlepage}

\tableofcontents
\newpage

% ---------------------------------------------------------------------------
% §1. Resumen (5 a 8 líneas)
% ---------------------------------------------------------------------------
\begin{abstract}
El presente trabajo analiza en qué medida el resultado de una carrera de
Fórmula 1 depende del rendimiento del auto o del talento del piloto. Se
utilizan los resultados de carrera y de clasificación del dataset
\emph{Formula 1 World Championship}, restringidos a las temporadas
2010--2024. Como los dos pilotos de una escudería disponen del mismo auto, la
diferencia entre compañeros de equipo se atribuye al piloto y la diferencia
entre escuderías al auto. Con esa idea se descompone, carrera por carrera, la
variación de las posiciones en un componente de escudería y uno de piloto, y
se estudia su evolución por temporada.
% TODO: agregar una o dos oraciones con los hallazgos principales cuando la
% sección de resultados esté cerrada.
\end{abstract}

% ---------------------------------------------------------------------------
% §2. Introducción (media página)
% ---------------------------------------------------------------------------
\section{Introducción}

\subsection{Motivación}
En la Fórmula 1 el resultado de cada carrera combina dos factores difíciles
de separar: el auto que construye la escudería y el piloto que lo maneja. Los
equipos invierten cientos de millones de dólares por año en el desarrollo del
auto y, al mismo tiempo, pagan salarios muy altos para contratar a los pilotos
más destacados. Estimar cuánto del resultado explica cada factor permite
evaluar dónde rinde más cada peso invertido, y también qué tan justo es
comparar a campeones que corrieron con autos de muy distinto nivel.

El período 2010--2024 resulta adecuado para estudiarlo porque contiene tres
reglamentos técnicos distintos (motores V8 hasta 2013, motores turbo híbridos
desde 2014 y autos de efecto suelo desde 2022) y la introducción del límite
presupuestario en 2021. Esos cambios permiten observar si el balance entre
auto y piloto se modifica cuando cambian las reglas que determinan cuán
parejos son los autos.

\subsection{Pregunta de investigación}
La pregunta que guía el trabajo es: \emph{¿qué pesa más en el resultado de una
carrera de Fórmula 1 entre 2010 y 2024, el rendimiento del auto o el talento
del piloto?} Para responderla se compara a cada piloto con su compañero de
equipo, que corre con el mismo auto en la misma carrera.

\subsection{Objetivos}
El objetivo general del trabajo es estimar qué proporción del resultado de
una carrera de Fórmula 1 se explica por el rendimiento del auto y qué
proporción por el desempeño del piloto en las temporadas 2010--2024, y
describir cómo cambió ese balance a lo largo del período.

Los objetivos particulares son:
\begin{enumerate}
    \item Construir, a partir de las tablas del dataset, una base de
    comparaciones entre compañeros de equipo (misma escudería, misma
    carrera), de modo que la diferencia de resultado entre ambos pilotos
    aísle el efecto del piloto manteniendo constante el auto.
    \item Descomponer la variación de las posiciones dentro de cada carrera en
    un componente atribuible a la escudería y otro atribuible al piloto, tanto
    en clasificación como en carrera, y compararlo con el valor esperable si
    la escudería no tuviera ningún efecto.
    \item Analizar la evolución del peso relativo de cada factor por
    temporada:
    \begin{enumerate}
        \item comparar las eras reglamentarias del período: motores V8
        (2010--2013), turbo híbridos (2014--2021) y efecto suelo
        (2022--2024);
        \item evaluar si a partir del límite presupuestario de 2021 los autos
        quedaron más parejos y aumentó el peso del piloto.
    \end{enumerate}
    \item Identificar a los pilotos que superan de forma sistemática a sus
    compañeros a lo largo de distintos equipos, como evidencia de un efecto
    piloto que se mantiene al cambiar de auto.
\end{enumerate}

\subsection{Dataset y fuente}
Se utiliza el dataset \emph{Formula 1 World Championship (1950--2024)},
publicado en Kaggle por Rohan Rao a partir de la base Ergast
\cite{kaggle}. Está compuesto por un conjunto de tablas relacionadas
(archivos CSV) que registran carreras, resultados, clasificaciones, pilotos,
escuderías y circuitos del campeonato mundial desde 1950. El análisis se
restringe a las temporadas 2010--2024: 305 carreras, 80 pilotos y 23
escuderías.

% ---------------------------------------------------------------------------
% §3. Datos y metodología (media página)
% ---------------------------------------------------------------------------
\section{Datos y metodología}

\subsection{Variables y estructura}
Se utilizan seis tablas del dataset: \texttt{races} (carreras, con año y
fecha), \texttt{results} (resultado de cada piloto en cada carrera),
\texttt{qualifying} (posición de clasificación), \texttt{drivers},
\texttt{constructors} y \texttt{status} (motivo por el que cada auto terminó
o abandonó). Las tablas se unen mediante sus claves \texttt{raceId},
\texttt{driverId}, \texttt{constructorId} y \texttt{statusId}. La tabla
principal resultante tiene 6.436 filas, una por piloto y carrera, y 25
columnas. No contiene filas duplicadas ni pilotos repetidos dentro de una
misma carrera.

Las variables de resultado son dos. En carrera se usa
\texttt{positionOrder}, que ordena a todos los autos, incluidos los que
abandonaron. En clasificación se usa \texttt{position}, la posición de salida
obtenida en la sesión de clasificación, que mide el ritmo del auto y del
piloto sin la interferencia de la estrategia, los choques o las roturas.

\subsection{Transformaciones}
Sobre la tabla principal se realizaron tres transformaciones:
\begin{enumerate}
    \item \textbf{Tratamiento de faltantes y tipos.} El dataset no codifica
    los faltantes como celdas vacías sino con el texto \texttt{\textbackslash
    N}, por lo que columnas numéricas como \texttt{position} se leían como
    texto. Se reemplazó ese valor por un nulo y se convirtieron las columnas a
    tipo numérico. Los 1.039 faltantes de \texttt{position} corresponden a
    autos que no terminaron la carrera, no a errores de carga.
    \item \textbf{Exclusión de abandonos en carrera.} Se creó una variable que
    indica si el auto terminó la carrera, aunque sea con vueltas de atraso.
    El 82\% de los resultados cumple esa condición. El resto se excluye del
    análisis de carrera, porque un abandono por falla mecánica castiga al
    piloto por un problema del auto, y uno por choque mezcla ambos factores.
    La clasificación no tiene ese problema y se usa como medida
    complementaria.
    \item \textbf{Comparación contra el compañero.} Para cada carrera y
    escudería se calculó la posición promedio de sus dos autos, que
    representa el nivel del auto en ese fin de semana, y la diferencia de
    cada piloto contra ese promedio, que representa el aporte del piloto. Se
    conservaron solo las escuderías con exactamente dos autos comparables en
    la carrera: 3.194 pares de compañeros en clasificación y 2.209 en carrera.
\end{enumerate}

\subsection{Método de análisis}
Para cada carrera, la variación total de las posiciones alrededor de la media
de la carrera se descompone en dos partes, con la misma lógica que un análisis
de varianza (ANOVA) de un factor:
\begin{equation}
    \sum_{i} (y_i - \bar{y})^2 =
    \underbrace{\sum_{i} (\bar{y}_{e(i)} - \bar{y})^2}_{\text{auto}} +
    \underbrace{\sum_{i} (y_i - \bar{y}_{e(i)})^2}_{\text{piloto}},
\end{equation}
donde $y_i$ es la posición del piloto $i$, $\bar{y}$ la posición media de la
carrera y $\bar{y}_{e(i)}$ la posición media de su escudería. El cociente
entre el componente del auto y el total indica qué porcentaje de la variación
explica la escudería. Ambas partes se suman por temporada con
\texttt{groupby} para obtener la evolución anual. Como referencia, si la
escudería no tuviera ningún efecto y los pilotos se repartieran al azar, con
10 equipos de dos autos el componente del auto explicaría igualmente cerca
del 47\% (9 de los 19 grados de libertad).

Como medida complementaria del talento se calcula, para cada piloto con al
menos 50 clasificaciones comparables, el porcentaje de duelos ganados a su
compañero y la diferencia promedio de posición contra él.

\subsection{Herramientas}
El análisis se realizó en Python sobre Google Colab, con las librerías
\texttt{pandas} y \texttt{numpy} para la manipulación de datos y
\texttt{matplotlib} para los gráficos. Los datos se descargan con
\texttt{kagglehub}; el notebook incluye un respaldo que obtiene los mismos
archivos desde un espejo público con el mismo esquema. El informe se
redactó en \LaTeX{} (Overleaf). El código y los gráficos están disponibles en
el repositorio del grupo.

% ---------------------------------------------------------------------------
% §4. Resultados (1 a 2 páginas)
% ---------------------------------------------------------------------------
\section{Resultados}
% TODO: un párrafo de interpretación por figura o tabla, que explique qué
% significa en relación con la pregunta (no solo qué se ve).

\subsection{Diferencias entre escuderías y entre compañeros}
% TODO: tabla de resumen con groupby (por ejemplo, dispersión de resultados
% entre equipos vs. dentro de cada equipo, por era).
% \begin{table}[H]
%     \centering
%     \caption{...}
%     \label{tab:dispersion}
%     \begin{tabular}{lrr}
%         \toprule
%         Era & Brecha entre equipos & Brecha entre compañeros \\
%         \midrule
%         ... & ... & ... \\
%         \bottomrule
%     \end{tabular}
% \end{table}

\subsection{Evolución por eras}
% TODO: gráfico de evolución temporal (líneas) del peso relativo del auto.
% \begin{figure}[H]
%     \centering
%     \includegraphics[width=0.85\textwidth]{evolucion_eras.png}
%     \caption{...}
%     \label{fig:evolucion}
% \end{figure}

\subsection{Comparación entre categorías}
% TODO: gráfico de barras (por ejemplo, porcentaje de duelos ganados por
% piloto, o resultados promedio por escudería dominante).

\subsection{Relación entre variables}
% TODO: dispersión entre rendimiento del auto (posición media del equipo) y
% resultado del piloto, o entre grilla de partida y posición final.

% ---------------------------------------------------------------------------
% §5. Conclusiones (media página)
% ---------------------------------------------------------------------------
\section{Conclusiones}
% TODO: responder la pregunta de forma directa, retomando cada objetivo
% particular.

\subsection{Limitaciones}
% TODO. Algunas a considerar:
%  - El período 2010--2024 no permite generalizar a las décadas anteriores
%    del campeonato, en las que la tecnología y los reglamentos eran otros.
%  - Hay pocos equipos por temporada (10 a 12), por lo que los resultados
%    por era se apoyan en pocas observaciones de escudería.
%  - Los reemplazos a mitad de temporada generan duelos con pocas carreras,
%    con menos peso estadístico que los duelos de temporada completa.
%  - La comparación entre compañeros mide al piloto respecto de otro piloto,
%    no en términos absolutos.

\subsection{Posibles extensiones}
% TODO: por ejemplo, un modelo de regresión con efectos fijos de piloto y de
% escudería-temporada, o incorporar los tiempos de vuelta (lap_times).

% ---------------------------------------------------------------------------
% §6. Reflexión metodológica (un cuarto de página)
% ---------------------------------------------------------------------------
\section{Reflexión metodológica}
% TODO: qué herramientas de IA se usaron y para qué (código, redacción,
% estructura del informe), qué se verificó manualmente, cómo fue la revisión
% intragrupal y, si se hizo, qué comentarios dejó la validación externa y
% cómo se incorporaron. Cerrar con uno o dos aprendizajes.

% ---------------------------------------------------------------------------
% §7. Bibliografía (opcional)
% ---------------------------------------------------------------------------
\begin{thebibliography}{9}
    \bibitem{kaggle}
    Rao, R. \emph{Formula 1 World Championship (1950--2024)}.
    Kaggle. \url{https://www.kaggle.com/datasets/rohanrao/formula-1-world-championship-1950-2020}
\end{thebibliography}

\end{document}
